\begin{blocksection}
\question Consider tree-recursive Fibonacci for non-negative integers.
\begin{lstlisting}
def fib(n):
    if n < 2:
        return n
    return fib(n - 1) + fib(n - 2)
\end{lstlisting}
\begin{parts}
\part Draw the calls made by \lstinline{fib(4)}. Which inputs occur more
than once, and what does the original call return?
\begin{solution}[1.4in]
\begin{verbatim}
fib(4)
  fib(3)
    fib(2)
      fib(1)
      fib(0)
    fib(1)
  fib(2)
    fib(1)
    fib(0)
\end{verbatim}
Inputs 2, 1, and 0 repeat. The result is 3.
\end{solution}
\part Suppose every call checks the same cache and saves its result.
What should the keys and values be? When the second request for
\lstinline{fib(2)} occurs, what work is avoided?
\begin{solution}[0.8in]
Keys are input integers; values are their Fibonacci numbers. The first
request stores 1 under key 2. The second returns that result without
calling \lstinline{fib(1)} or \lstinline{fib(0)} again. Each distinct
Fibonacci result is computed once, even if requested repeatedly.
\end{solution}
\end{parts}
\end{blocksection}
